% Part: first-order-logic
% Chapter: sequent-calculus
% Section: proving-things-quant

\documentclass[../../../include/open-logic-section]{subfiles}

\begin{document}

\olfileid{fol}{seq}{prq}

\olsection{పరిమాణీకరణలతో \usetoken{P}{derivation}}

\begin{ex}
$\lexists[x][\lnot !A(x)] \Sequent \lnot \lforall[x][!A(x)]$ అనే
సీక్వెంట్‌కు ఒక $\Log{LK}$-\tetoken{వ్యుత్పత్తి}{!!{derivation}} ఇవ్వండి.

పరిమాణీకరణలతో పనిచేసేటప్పుడు ఐగెన్ చరరాశి షరతును ఉల్లంఘించకుండా
చూసుకోవాలి. కొన్నిసార్లు దానికి కొన్ని నిగమనాలను అమలు చేసే క్రమాన్ని
మార్చి ప్రయత్నించాల్సి వస్తుంది. సాధారణంగా ఐగెన్ చరరాశి షరతుకు లోబడే
నియమాలను ముందుగా చూసుకోవడం సహాయపడుతుంది (పూర్తయిన నిరూపణలో అవి
కిందివైపు ఉంటాయి). ముందుచూపుతో ప్రారంభ సీక్వెంట్ ఎలా ఉండవచ్చో
ఊహించడానికి ప్రయత్నించడం కూడా మంచిది. మన సందర్భంలో అది
$!A(a) \Sequent !A(a)$ లాంటిది అయి ఉండాలి. అంటే పరిమాణీకరణ నియమాలను
“వెనక్కు” వర్తింపజేసేటప్పుడు $\lforall$, $\lexists$ నియమాలు రెండింటికీ
ఒకే పదాన్ని---దానిని $a$ అందాం---ఎంచుకోవాలి. ఒక్కో నియమానికి వేర్వేరు
పదాలను ఎంచుకుంటే చివరికి $!A(a) \Sequent !A(b)$ వంటిది వస్తుంది;
అది నిగమించదగినది కాదని స్పష్టం.

ఎప్పటిలాగే ఇలా రాయడం మొదలుపెడతాం:
\begin{prooftree}
\AxiomC{}
\UnaryInf$\lexists[x][\lnot !A(x)] \fCenter \lnot \lforall[x][!A(x)]$
\end{prooftree}
\LeftR{\exists} లేదా \RightR{\lnot} నియమాన్ని వర్తింపజేయవచ్చు.
\LeftR{\exists} నియమం ఐగెన్ చరరాశి షరతుకు లోబడుతుంది కనుక దాన్ని
ఆలస్యం చేయకుండా ముందే చూసుకోవడం మంచిది; అందువల్ల మొదట దానినే
వర్తింపజేద్దాం.
\begin{prooftree}
\AxiomC{}
\UnaryInf$ \lnot !A(a) \fCenter \lnot \lforall[x][!A(x)]$
\RightLabel{\LeftR{\lexists}}
\UnaryInf$ \lexists[x][\lnot !A(x)] \fCenter \lnot \lforall[x][!A(x)]$
\end{prooftree}
\LeftR{\lnot}, \RightR{\lnot} నియమాలను వెనక్కు వర్తింపజేస్తే ఇలా వస్తుంది:
\begin{prooftree}
\AxiomC{}
\UnaryInf$\lforall[x][!A(x)] \fCenter !A(a)$
\RightLabel{\LeftR{\lnot}}
\UnaryInf$\lnot !A(a), \lforall[x][!A(x)] \fCenter $
\RightLabel{\LeftR{\Exchange}}
\UnaryInf$\lforall[x][!A(x)], \lnot !A(a) \fCenter $
\RightLabel{\RightR{\lnot}}
\UnaryInf$ \lnot !A(a) \fCenter \lnot \lforall[x] !A(x)$
\RightLabel{\LeftR{\lexists}}
\UnaryInf$ \lexists[x] \lnot !A(x) \fCenter \lnot \lforall[x] !A(x)$
\end{prooftree}
ఈ దశలో \LeftR{\forall} నియమాన్ని వర్తింపజేయడమే మనకున్న ఏకైక
ఎంపిక. ఈ నియమం ఐగెన్ చరరాశి పరిమితికి లోబడదు కనుక ఇబ్బంది లేదు.
ప్రారంభ సీక్వెంట్‌ను ($!A(a) \Sequent !A(a)$ రూపంలో ఉన్నదాన్ని)
పొందడానికి ప్రయత్నిస్తున్నామని గుర్తుంచుకోండి. అందువల్ల నియమాన్ని
వర్తింపజేసేటప్పుడు $!A$కు ఆర్గ్యుమెంటుగా $a$నే ఎంచుకోవాలి.
\begin{prooftree}
\Axiom$!A(a) \fCenter !A(a)$
\RightLabel{\LeftR{\lforall}}
\UnaryInf$\lforall[x][!A(x)] \fCenter !A(a)$
\RightLabel{\LeftR{\lnot}}
\UnaryInf$\lnot !A(a), \lforall[x][!A(x)] \fCenter $
\RightLabel{\LeftR{\Exchange}}
\UnaryInf$\lforall[x][!A(x)], \lnot !A(a) \fCenter $
\RightLabel{\RightR{\lnot}}
\UnaryInf$ \lnot !A(a) \fCenter \lnot \lforall[x][!A(x)]$
\RightLabel{\LeftR{\lexists}}
\UnaryInf$ \lexists[x][ \lnot !A(x)] \fCenter \lnot \lforall[x][!A(x)]$
\end{prooftree}
ముఖ్యంగా పరిమాణీకరణలతో పనిచేసేటప్పుడు, ఈ దశలో ఐగెన్ చరరాశి షరతు
ఉల్లంఘించబడలేదని మరొకసారి తనిఖీ చేయాలి. మనం వర్తింపజేసిన నియమాల్లో
ఐగెన్ చరరాశి షరతుకు లోబడింది \LeftR{\exists} మాత్రమే; దాని ఐగెన్
చరరాశి~$a$ దాని కింది సీక్వెంట్‌లో (అదే అంత్య సీక్వెంట్‌లో) కనిపించదు.
కాబట్టి ఇది సరైన \tetoken{వ్యుత్పత్తి}{!!{derivation}}.
\end{ex}

\begin{prob}
కింది సీక్వెంట్లకు \tetoken{వ్యుత్పత్తులు}{!!{derivation}s} ఇవ్వండి:
\begin{enumerate}
\item $\Sequent (\lforall[x][!A(x)] \land \lforall[y][!B(y)]) \lif
\lforall[z][(!A(z) \land !B(z))]$.
\item $\Sequent (\lexists[x][!A(x)] \lor \lexists[y][!B(y)]) \lif
\lexists[z][(!A(z) \lor !B(z))]$.
\item $\lforall[x][(!A(x) \lif !B)] \Sequent \lexists[y][!A(y)] \lif !B$.
\item $\lforall[x][\lnot !A(x)] \Sequent \lnot\lexists[x][!A(x)]$.
\item $\Sequent \lnot\lexists[x][!A(x)] \lif \lforall[x][\lnot !A(x)]$.
\item $\Sequent \lnot\lexists[x][\lforall[y][((!A(x,y) \lif \lnot
!A(y,y)) \land (\lnot !A(y,y) \lif !A(x,y)))]]$.
\end{enumerate}
\end{prob}

\begin{prob}
కింది సీక్వెంట్లకు \tetoken{వ్యుత్పత్తులు}{!!{derivation}s} ఇవ్వండి:
\begin{enumerate}
\item $\Sequent \lnot\lforall[x][!A(x)] \lif \lexists[x][\lnot!A(x)]$.
\item $(\lforall[x][!A(x)] \lif !B) \Sequent \lexists[y][(!A(y) \lif !B)]$.
\item $\Sequent \lexists[x][(!A(x) \lif \lforall[y][!A(y)])]$.
\end{enumerate}
(వీటన్నింటికీ $\RightR{\Contraction}$ నియమం అవసరం.)
\end{prob}

\end{document}
